\section{Book Shop}
\label{sec:cses-book-shop}

\problemstatement{You have $x$ money and there are $n$ books, each with a
price and a page count. You may buy each book at most once. Maximise the
total number of pages you can buy within your budget.}

\begin{patternbox}
\emph{``Each item at most once, maximise value under a capacity''} is the
classic \textbf{0/1 knapsack}. Price is weight, pages are value, $x$ is
capacity. Let $\textit{dp}[b]$ be the maximum pages achievable with
budget exactly $\le b$.
\end{patternbox}

\subsection*{State, transition, and the reverse inner loop}

Process books one at a time. $\textit{dp}[b]$ = best pages using books
considered so far with budget $b$. For a book with price $p$ and pages
$s$, either skip it ($\textit{dp}[b]$ unchanged) or buy it
($\textit{dp}[b-p]+s$). The decisive implementation detail: when using a
1D array, iterate $b$ \textbf{downward} from $x$ to $p$, so that
$\textit{dp}[b-p]$ still refers to a state \emph{without} the current
book -- guaranteeing each book is used at most once.
\[
  \textit{dp}[b]=\max\big(\textit{dp}[b],\ \textit{dp}[b-p]+s\big).
\]

\begin{ideabox}[Downward Loop Enforces ``At Most Once'']
In 0/1 knapsack the inner budget loop must go high-to-low. A low-to-high
loop would let $\textit{dp}[b-p]$ already include the current book,
effectively allowing it to be bought repeatedly -- which is the unbounded
knapsack (Coin Combinations II). The loop direction is the whole
difference.
\end{ideabox}

\subsection*{Algorithm}

\begin{enumerate}
  \item $\textit{dp}[b]=0$ for all $b$.
  \item For each book $(p,s)$: for $b=x$ down to $p$,
        $\textit{dp}[b]=\max(\textit{dp}[b],\textit{dp}[b-p]+s)$.
  \item Answer $\textit{dp}[x]$.
\end{enumerate}

\subsection*{Dry run}

\dryrun{Budget $x=10$; books $(\text{price},\text{pages})=(4,5),(8,12),
(5,8),(3,1)$.}

\begin{itemize}
  \item After book $(4,5)$: $\textit{dp}[b]=5$ for $b\ge4$.
  \item After $(5,8)$: $\textit{dp}[9]=\max(5,\textit{dp}[4]+8)=13$
        (buy the $4$ and the $5$).
  \item Final $\textit{dp}[10]=13$ -- books priced $4$ and $5$ give
        $5+8=13$ pages within budget $10$.
\end{itemize}

\subsection*{C++ implementation}

\begin{lstlisting}
#include <bits/stdc++.h>
using namespace std;

int main() {
    int n, x;
    cin >> n >> x;
    vector<int> price(n), pages(n);
    for (int& p : price) cin >> p;
    for (int& s : pages) cin >> s;

    vector<int> dp(x + 1, 0);
    for (int i = 0; i < n; i++)
        for (int b = x; b >= price[i]; b--)        // downward: 0/1
            dp[b] = max(dp[b], dp[b - price[i]] + pages[i]);

    cout << dp[x] << "\n";
    return 0;
}
\end{lstlisting}

\begin{complexitybox}
\textbf{Time Complexity.} $O(nx)$ -- each book sweeps the budget once.
\textbf{Space Complexity.} $O(x)$ with the rolling 1D array.
\end{complexitybox}

\begin{takeawaybox}
Book Shop is 0/1 knapsack in disguise; the only thing to get right is the
downward inner loop that forbids reusing an item. Contrast with the coin
problems' upward loop (items reusable). Loop direction is the recurring
knapsack tell.
\end{takeawaybox}
